% PROJECT SUMMARY / ABSTRACT -- NIH R01. Limit: 30 lines of text.
% Written for a broad scientific audience; no proprietary or confidential info.
\documentclass[11pt]{article}
\usepackage{definitions}

\makeatletter

\begin{document}

\linenumbers % Start numbering lines for the entire document

\textbf{PROJECT SUMMARY}

About 1 in 2,000 US newborns has single-ventricle heart disease (SVHD), treated by rerouting the systemic venous blood directly to the lungs (Fontan circulation). The Fontan circulation is palliative, as many patients progress to heart or multi-organ failure. However, some patients have a borderline-size ventricle: underdeveloped, yet with the potential to serve as a working ventricle. While such patients historically undergo a single-ventricle palliation, nowadays, leading clinical centers are developing complex biventricular (BiV) repair programs. In some patients, the small ventricle can substantially grow after surgically redirecting more blood into this ventricle (the so-called ventricular recruitment (VR)), and some hearts with a borderline-sized ventricle can be effective when the two ventricles are fully septated and the systemic and pulmonary circulations separated (BiV repair). This frees a child from the Fontan physiology for life. It is estimated that up to a third of Fontan palliations could be avoided. However, the risks associated with these complex procedures need to be assessed. Namely, for the BiV repair, it is the failure of the ventricle’s ability to sustain circulation; and for VR, the risk of insufficient growth in the next 6-15 months (leading to eventual re-entry into the single-ventricle palliation). The decision about the optimal clinical is currently based mainly on using heuristic criteria, e.g., end-diastolic volume or pressure of the borderline-sized ventricle, with a limited weight on functional indicators. Rather than analyzing sources of clinical data separately and seeking optimal cut-off values (flawed by inherent measurement inaccuracies and heterogeneity in patient population), we will synthesize the multimodality data – cardiac magnetic resonance imaging and pressure catheter – acquired prior to surgery, within a personalized biomechanical modeling framework, and create the so-called Digital Twins (DTs). Such a model-augmented assessment will provide a comprehensive analysis of the cardiovascular physiology prior to surgery, which will be in line with physiological and biophysical principles. Subsequently, the surgical procedure (VR or BiV repair) will be performed in silico using the DTs. The capability of short-term biomechanical models to predict the re-adjustment of cardiovascular physiology after BiV repair, and of long-term growth/remodeling (G/R) models to predict the outcome of VR, will be confronted with the outcomes of actual surgery. The study brings together investigative teams of clinical and biomechanical research, with a primary site located within clinical environment.  The overarching goal is to reduce the number of patients with SVHD requiring the Fontan palliation. The project will pave the way for DTs assisting in optimal clinical management for this patient group and for a subsequent paradigm shift in various decisions related to pediatric cardiology and other specialties.

\end{document}